\documentclass[10pt,letterpaper]{article}
\usepackage{technical-design-spec}
\setDocTitle{Formal Security Models and Assurance}
\setDocSubtitle{Protection state, lattice policies, noninterference and evidence}
\setDocVersion{1.0}
\setDocStatus{Independent educational companion}
\setDocOwner{latex-docs / security}
\setDocDate{2026-11-02}
\setDocSummary{Independent model comparison and formal-methods bridge keyed to Bishop Chapters 2--9 and 17--22.}
\begin{document}
\makedoctocpage
\begin{notebox}[Source and authority boundary]
\textbf{BOOK:} Matt Bishop, \emph{Computer Security: Art and Science}, Second Edition (Pearson, 2019). Book-aligned sections are \emph{original summaries and reading prompts}, not copied textbook content.
\textbf{EXTENSION:} Any contemporary framework mappings, implementation tasks, scoring rules, and operational recommendations are independently authored educational supplements, not the book's positions. Nothing here proves compliance.
\end{notebox}
\section{Model comparison and domain of validity}
\begin{longtable}{@{}p{0.20\linewidth}p{0.34\linewidth}p{0.35\linewidth}@{}}
\toprule Model & Property it expresses & Assumptions / limitations\\\midrule\endhead
Access control matrix & $M[s,o]$ defines rights of subject $s$ on object $o$ & The subject/object universe and transitions must be specified\\
HRU safety & Whether rights can leak into a specified matrix cell & The general safety problem is undecidable; bounded enumeration proves only the bounded case\\
Take--grant & Rights propagation through take and grant graph edges & Interpreted only under defined graph rewrite rules\\
Bell--LaPadula & Confidentiality lattice constraints & Trusted subjects, classification labels and declassification need explicit treatment\\
Biba & Integrity ordering and flow restrictions & Strict policy and relaxed variants make different trade-offs\\
Clark--Wilson & Well-formed transactions on constrained data items & Certification of transformation procedures and separation of duties are crucial\\
Chinese Wall & History-dependent conflict-of-interest restrictions & Conflict classes and access history must be accurate\\
RBAC & Role membership, permissions and separation constraints & Role explosion and context sensitivity are outside elementary models\\
Noninterference & High secrets do not alter low-observable behavior & Observable model and assumptions about timing/concurrency matter\\\bottomrule
\end{longtable}
\section{Access matrix transition system}
Let $S$ be subjects, $O$ objects, $R$ rights, and $M:S\times O\rightarrow\mathcal{P}(R)$. A protection state is $(S,O,M)$; an operation maps a state to a successor only when its guard holds. Specify \emph{creation}, \emph{grant}, \emph{revocation} and \emph{ownership} transitions separately. A bounded reachability search can exhibit a violation trace but cannot establish decidability of general HRU safety.
\section{Confidentiality and integrity label orders}
For security labels $L$ with order $\preceq$, specify whether higher values mean greater secrecy or greater integrity \emph{before} applying shorthand. In the standard simplified confidentiality lattice, prohibit a subject from reading a more confidential object and from writing into a less confidential object. In strict Biba integrity, reading from lower-integrity objects and writing to higher-integrity objects are constrained. These are \emph{different goals}; combining them requires explicit policy choices.
\section{Clark--Wilson and workflow correctness}
Define constrained data items (CDIs), unconstrained data items (UDIs), integrity verification procedures (IVPs) and transformation procedures (TPs). Certification and enforcement relationships must assign authorized subjects to certified TPs; an audit trail should record their use. A business approval workflow is not an instance of Clark--Wilson solely because it has two signatories.
\section{Noninterference and the two-run condition}
Let $h$ vary across high-security inputs, let $l$ be fixed low-security input, and let $\pi_L(P(h,l))$ be low-observable output. A simple deterministic noninterference test asks whether
\[
 \forall h_1,h_2,l:\quad \pi_L(P(h_1,l))=\pi_L(P(h_2,l)).
\]
This equation is a pedagogical abstraction. Termination, timing, concurrency and probabilistic channels require a stronger, specified observational model; two sample executions are not proof over all inputs.
\section{Assurance argument pattern}
\begin{enumerate}
\item Claim: identify precise asset and security property.
\item Context: define system boundary, environment, attacker and trusted base.
\item Argument: explain why verified mechanisms imply the claim under assumptions.
\item Evidence: provide formal proof obligations, tests, review records and configuration hashes.
\item Defeaters: list counterexamples, unmodeled channels and evidence expiration conditions.
\end{enumerate}
\section{Evaluation and historical boundaries}
Bishop's modeling tools are foundational concepts. Product certification, deployment-specific risk decisions and contemporary controls must be established separately against authoritative criteria. Consult the published ISO/IEC 15408 series where evaluation is relevant; do not invent Common Criteria certification for a training program.

\end{document}
